\bookchapter{Kapitel 14: Passwort-Manager im Detail}

\chapterintro{Im digitalen Zeitalter ist dein Passwort der Schlüssel zu deiner Haustür. Aber bei zu vielen, zu komplexen Schlüsseln: Was tun? Der Passwort-Manager ist dein „Universalverwalter“ für Schlüssel.}

\booksection{I. Warum brauchst du einen Passwort-Manager?}

Ein moderner Mensch hat im Schnitt Dutzende bis Hunderte Online-Konten. Nutzt du für alle dasselbe Passwort, sind nach einem Hack alle Konten gefährdet. Nutzt du für jedes Konto ein anderes, komplexes Passwort, kannst du dir keines merken.

Der Passwort-Manager löst dieses Dilemma: Du musst dir nur ein einziges Master-Passwort merken – er merkt sich alle anderen für dich.

\booksection{II. Wie Passwort-Manager funktionieren}

\begin{enumerate}
  \item \textbf{Verschlüsselte Speicherung}: Alle Passwörter werden hochverschlüsselt gespeichert – nur du kannst sie entschlüsseln.
  \item \textbf{Master-Passwort}: Du legst ein Master-Passwort fest, das den Tresor entsperren kann. Das Master-Passwort selbst wird nicht auf dem Server gespeichert.
  \item \textbf{Automatisches Ausfüllen}: Auf Login-Seiten von Websites und Apps füllt der Passwort-Manager Konto und Passwort automatisch aus.
  \item \textbf{Passwort-Generator}: Er kann zufällige, komplexe Passwörter erzeugen und so die Sicherheit erhöhen.
\end{enumerate}

\booksection{III. Kurze Einfuehrung in die Verschluesselung}

Das Verstaendnis der Verschluesselungsprinzipien von Passwortmanagern hilft dir einzuschaetzen, ob ein Anbieter vertrauenswuerdig ist. Im Folgenden werden Kernkonzepte in verstaendlicher Sprache erklaert.

\booksubsection{AES-256-Verschluesselung}

AES (Advanced Encryption Standard) ist ein vom NIST anerkannter symmetrischer Verschluesselungsalgorithmus, der von Banken, dem Militaer und Regierungen weit verbreitet ist. „256“ bezeichnet die Schluessellaenge von 256 Bit – ein Bruetforce-Angriff wuerde laenger dauern als das Alter des Universums. Massgebliche Passwortmanager (1Password, Bitwarden, KeePass u. a.) verschluesseln deine Passwortdatenbank mit AES-256.

\booksubsection{Zero-Knowledge-Architektur}

„Zero Knowledge“ bedeutet nicht, dass der Anbieter nichts speichert, sondern dass\textbf{Verschluesselung und Entschluesselung lokal auf deinem Geraet stattfinden; der Server speichert nur denverschluesselten Geheimtext}. Konkreter Ablauf:

\begin{enumerate}
  \item Du gibst das Master-Passwort im Client ein.
  \item Der Client leitet aus dem Master-Passwort den Verschluesselungsschluessel ab (dieser Schritt erfolgt auf deinem Smartphone/Computer).
  \item Nur der verschluesselte Geheimtext wird zum Server hochgeladen.
  \item Beim naechsten Login sendet der Server den Geheimtext; der Client entschlüsselt ihn erneut mit dem Master-Passwort.
\end{enumerate}

Wesentlicher Punkt:\textbf{Das Master-Passwort verlaesst niemals dein Geraet}, auch nicht auf dem Server. Selbst wenn der Anbieter gehackt wird, erhaelt der Angreifer nur unentschluesselbaren Geheimtext – es sei denn, er koennte dein Master-Passwort erraten.

\begin{tcbinfo}{Warum bedeutet „Zero Knowledge", dass auch der Anbieter deine Passwoerter nicht sehen kann?}
Weil der Entschluesselungsschluessel aus deinem Master-Passwort abgeleitet wird und dieses nur in deinem Kopf und lokal auf deinem Geraet existiert. Auf den Servern des Anbieters gibt es weder das Master-Passwort noch den Entschluesselungsschluessel – nur den Geheimtext. Deshalb erklaeren seriöse Anbieter wie Bitwarden und 1Password ausdruecklich: „Auch wir koennen deine Passwoerter nicht sehen“ – nicht weil sie nicht wollten, sondern weil es technisch unmoeglich ist.

\end{tcbinfo}

\booksubsection{Schluesselableitungsfunktionen: PBKDF2 und Argon2}

Das Master-Passwort kann nicht direkt als Verschluesselungsschluessel verwendet werden – es muss durch eine „Schluesselableitungsfunktion“ verarbeitet werden. Ihre Aufgabe ist es,\textbf{die Berechnung absichtlich zu verlangsamen}, um Bruetforce-Angriffe so teuer zu machen, dass sie unrentabel werden.

\begin{itemize}
  \item \textbf{PBKDF2}: Etablierter Standard, verwendet von iCloud-Schluesselbund, Bitwarden, LastPass u. a.; verlangsamt Angriffe durch wiederholte Iteration (meist ueber 100.000-mal) des Hashs.
  \item \textbf{Argon2}: Sieger des Password Hashing Competition 2015; moderner als PBKDF2. Neben „Rechenzeit" auch „Speicherverbrauch" – GPUs/ASICs koennen Angriffe nicht beschleunigen. Bitwarden, KeePassXC u. a. unterstuetzen es bereits.
\end{itemize}

Hoehere Iterationszahlen sind sicherer, aber das Entsperren wird langsamer. Serioese Anbieter veroeffentlichen ihre Ableitungsalgorithmen und Iterationszahlen zur Pruefung durch die Security-Community.

\booksubsection{Fallbeispiel: LastPass-Datenleck 2022}

Im August 2022 meldete der etablierte Passwortmanager LastPass einen Einbruch in seine Entwicklungsumgebung; im November wurde bekannt, dass Angreifer Backups der Passwortdatenbanken der Nutzer (verschluesselter Geheimtext) entwendet hatten. Die wesentlichen Lehren:

\begin{itemize}
  \item \textbf{Die Zero-Knowledge-Architektur hielt die Verteidigungslinie}: Da die Passwortdatenbank verschluesselt war, erhielt der Angreifer den Geheimtext, konnte ihn aber nicht entschluesseln – vorausgesetzt, das Master-Passwort des Nutzers war stark genug.
  \item \textbf{Nutzer mit schwachem Master-Passwort am gefaehrdetsten}: Angreifer koennen schwache Master-Passwoerter offline per Bruetforce knacken; Nutzer mit schwachen oder wiederverwendeten Master-Passwoertern sind real gefaehrdet.
  \item \textbf{Auch Metadaten wurden geleakt}: Passwortinhalte waren zwar verschluesselt, aber Metadaten wie URLs und Notizen waren moeglicherweise unverschluesselt und zeigen, welche Dienste der Nutzer verwendet.
\end{itemize}

\begin{tcbwarn}{Lehren aus dem LastPass-Vorfall}
Erstens: Das Master-Passwort muss lang und einzigartig sein (mindestens 16 Zeichen, mit Gross-/Kleinbuchstaben, Zahlen und Symbolen). Zweitens: Bevorzuge Anbieter, die Argon2id und hohe Iterationszahlen unterstuetzen. Drittens: Fuer hochwertige Konten (Bank, E-Mail) zwingend eine unabhaengige Zwei-Faktor-Authentifizierung aktivieren – nicht nur auf den Passwortmanager vertrauen.

\end{tcbwarn}

\booksection{IV. Empfohlene Passwortmanager}

\booksubsection{Systemeigene:}

\begin{itemize}
  \item \textbf{iCloud-Schlüsselbund}: Erste Wahl für Apple-Nutzer. Nahtlos in Safari und Apps integriert, kostenlos, sicher. Nachteil: Plattformübergreifende Nutzung unpraktisch.
  \item \textbf{Google Passwortmanager}: Bequem für Android-Nutzer, in Chrome und Android integriert, kostenlos.
  \item \textbf{Huawei Passwort-Tresor}: In HarmonyOS integriert, für Huawei-Nutzer.
\end{itemize}

\booksubsection{Professionelle Drittanbieter:}

\begin{itemize}
  \item \textbf{Bitwarden}: Open Source, kostenlos, funktionsreich, unterstützt alle Plattformen. Für datenschutzbewusste Nutzer wärmstens empfohlen.
  \item \textbf{1Password}: Schöne Oberfläche, reich an Funktionen, geeignet für Familie und Team, kostenpflichtig.
  \item \textbf{LastPass}: Bekannter Passwortmanager, die kostenlose Version ist eingeschränkt.
  \item \textbf{KeePass}: Vollständig lokale Speicherung, Open Source und kostenlos – für Power-User.
\end{itemize}

\booksection{V. Realitaetsnahe Bewertung der Passwortstaerke}

Passwortmanager zeigen meist einen „Passwortstaerke“-Indikator (rot/gelb/gruen), der jedoch haeufig nur Laenge und Zeichenvielfalt betrachtet – unzureichend. Um die echte Sicherheit einzuschaetzen, muessen „Entropie” und „bereits geleakt?“ beruecksichtigt werden.

\booksubsection{Entropie: Die wissenschaftliche Messung der Passwortstaerke}

„Entropie“ misst die Unsicherheit eines Passworts in Bit. Hoehere Entropie bedeutet mehr Versuche fuer einen Bruetforce-Angriff.

\begin{itemize}
  \item \textbf{"123456"}: Sehr geringe Entropie, in Sekunden geknackt.
  \item \textbf{"P@ssw0rd"}: Scheint komplex, steht aber auf jeder Woerterbuch-Angriffsliste; Entropie weiterhin niedrig.
  \item \textbf{"correct horse battery staple"}: Vier zufaellige Woerter, ca. 44 Bit Entropie; schwerer zu knacken als „P@ssw0rd".
  \item \textbf{16 zufaellige Zeichen (Gross-/Kleinbuchstaben + Zahlen + Symbole)}: ca. 100 Bit Entropie; mit heutiger Rechenleistung nicht per Bruetforce knackbar.
\end{itemize}

Fazit:\textbf{„Laenge + Zufaelligkeit“ ist wichtiger als „komplexe Symbole”}. Genau deshalb sind die zufaelligen Langpasswoerter eines Passwortmanagers wirklich sicher.

\booksubsection{HaveIBeenPwned: Pruefe, ob dein Passwort bereits geleakt wurde}

Vom Security-Forscher Troy Hunt betreute\href{https://haveibeenpwned.com}{haveibeenpwned.com}Datenbank mit Konten und Passwoertern aus grossen Lecks. Du kannst:

\begin{enumerate}
  \item E-Mail-Adresse eingeben und pruefen, ob dein Konto in einem Leak aufgetaucht ist.
  \item Mit der Pwned-Passwords-Funktion pruefen, ob ein Passwort in einer bekannten Leak-Datenbank enthalten ist (Passwortmanager und Browser haben dies bereits integriert).
\end{enumerate}

Wenn ein Passwort in einer Leak-Datenbank auftaucht,\textbf{muss es unabhaengig von der scheinbaren Komplexitaet sofort geaendert werden}, denn Woerterbuchangriffe probieren diese bekannten Passwoerter zuerst.

\begin{tcbinfo}{Prinzip der k-Anonymitaets-Abfrage}
Sorge, dass die Eingabe bei HaveIBeenPwned zu einem erneuten Leak fuehrt? Keine Sorge. Moderne Abfragen nutzen das „k-Anonymitaets“-Protokoll: Der Browser sendet nur die ersten 5 Zeichen des Passwort-Hashs an den Server; dieser liefert alle Hashes mit diesem Praefix zurueck; der Abgleich erfolgt lokal. Der Server sieht niemals deinen vollstaendigen Passwort-Hash. Die „Leak-Erkennung” von 1Password und Bitwarden funktioniert so.

\end{tcbinfo}

\booksubsection{Integrierte Ueberwachungsfunktionen des Passwortmanagers}

\begin{itemize}
  \item \textbf{1Password Watchtower}: Ueberwacht in Echtzeit schwache, duplizierte, geleakte und abgelaufene Passwoerter sowie Seiten ohne 2FA; Anzeige in den Stufen rot/orange/gelb.
  \item \textbf{Bitwarden Vault Health}: Die Gratis-Version zeigt schwache und duplizierte Passwoerter; die Premium-Version erkennt geleakte Passwoerter (basierend auf HaveIBeenPwned).
  \item \textbf{iCloud-Schluesselbund „Sicherheitsempfehlungen"}: iOS Einstellungen Passwoerter Sicherheitsempfehlungen; zeigt, welche Passwoerter in bekannten Leaks auftauchen oder doppelt verwendet werden.
  \item \textbf{Google Passwortmanager „Passwortpruefung"}: In Chrome integriert; Ein-Klick-Pruefung auf geleakte, schwache und duplizierte Passwoerter.
\end{itemize}

\begin{tcbpractice}{Jetzt pruefen: Ist dein Passwort sicher?}
\begin{enumerate}
  \item Oeffne haveibeenpwned.com, gib deine haeufig verwendete E-Mail-Adresse ein und pruefe, in welchen Leaks sie auftaucht.
  \item iPhone-Nutzer: Einstellungen Passwoerter Sicherheitsempfehlungen; jeden roten Eintrag abarbeiten.
  \item Chrome-Nutzer: Einstellungen Passwoerter Passwoerter pruefen; nach Anleitung aendern.
  \item 1Password-/Bitwarden-Nutzer: Oeffne Watchtower/Health Reports; zuerst „geleakt" und „dupliziert" behandeln.
\end{enumerate}

\end{tcbpractice}

\booksection{VI. Passkey: Die Zukunft des passwortlosen Logins}

2024-2026 durchlaeuft die Passworttechnologie eine Revolution –\textbf{Passkey}koennte traditionelle Passwoerter vollstaendig abloesen. Apple, Google und Microsoft treiben es gemeinsam voran; unterstuetzt von der FIDO Alliance und dem W3C – die groesste Login-Revolution seit Jahrzehnten.

\booksubsection{Was ist ein Passkey?}

Passkey basiert auf dem FIDO2-/WebAuthn-Standard und verwendet\textbf{ein Schluesselpaar auf dem Geraet}anstelle eines Passworts. Das Prinzip vereinfacht:

\begin{itemize}
  \item Dein Smartphone erzeugt ein Schluesselpaar:\textbf{Public Key}(bei der Website registriert) +\textbf{Private Key}(bleibt immer im Sicherheitschip deines Geraets).
  \item Beim Login sendet die Website eine „Challenge"; dein Smartphone signiert sie mit dem Private Key, die Website verifiziert mit dem Public Key.
  \item Vor der Signatur musst du dich per Biometrie (Face ID / Fingerabdruck) oder Geraete-PIN authentifizieren.
  \item Waehrend des gesamten Vorgangs\textbf{wird kein Passwort ueber das Netzwerk uebertragen}.
\end{itemize}

\booksubsection{Passkey vs. traditionelle Passwoerter}

\begin{center}
\begin{tabularx}{\linewidth}{|X|X|X|}
\hline
\textbf{Dimension} & \textbf{Traditionelles Passwort} & \textbf{Passkey} \tabularnewline
\hline
Phishing-Schutz & Schlecht; Phishing-Seiten koennen Passwoerter stehlen & Phishing-resistent; Private Key signiert nur fuer echte Domains \tabularnewline
Auswirkung eines Server-Lecks & Passwort-Hashs koennten geknackt werden & Server speichert nur Public Key; Leck ohne Auswirkung \tabularnewline
Komfort & Muss gemerkt/eingegeben/verwaltet werden & Face ID/Fingerabdruck – Login in einer Sekunde \tabularnewline
Replay-Angriff & Kann per MITM abgefangen und wiedergegeben werden & Jede Signatur ist einmalig; nicht wiederverwendbar \tabularnewline
\hline
\end{tabularx}
\end{center}

\booksubsection{Plattform-Unterstuetzung}

\begin{itemize}
  \item \textbf{Apple iCloud-Schluesselbund Passkey}: iOS 16+ / macOS Ventura+ nativ; Passkeys synchronisieren zwischen allen Geraeten mit derselben Apple ID; auf neuen Geraeten keine Neuerstellung erforderlich.
  \item \textbf{Google Passwortmanager Passkey}: Android + Chrome nativ; Passkeys werden im Google-Konto gespeichert und synchronisieren zwischen Android-Geraeten; seit 2023 unterstuetzt das persoenliche Google-Konto Passkey-Login.
  \item \textbf{Microsoft}: Windows 11 / Edge unterstuetzen Passkey; Microsoft-Konto, Outlook, Office und weitere Dienste eroeffnen nach und nach Passkey-Login; in Microsoft Authenticator verwaltbar.
  \item \textbf{1Password / Bitwarden}: Unterstuetzt Erstellung, Speicherung und plattformuebergreifende Synchronisation von Passkeys; unabhaengig von einem einzigen Oekosystem.
\end{itemize}

\begin{tcbinfo}{Plattformuebergreifende Synchronisation als Durchbruch}
Frueher war die groesste Schwachstelle von Passkeys die „Bindung an ein einzelnes Geraet“ – bei Geraetewechsel musste man neu erstellen. Seit 2023 haben Apple, Google und Microsoft die Cloud-Synchronisation von Passkeys umgesetzt: iCloud-Schluesselbund synchronisiert zwischen Apple-Geraeten, Google Passwortmanager zwischen Android-Geraeten. 1Password und weitere Drittanbieter ermoeglichen die plattformuebergreifende Synchronisation. Damit ist Passkey praktisch einsatzbereit.

\end{tcbinfo}

\booksubsection{So erstellst und verwendest du Passkeys auf dem Smartphone}

\textbf{iOS (iPhone) Pfad:}

\begin{enumerate}
  \item Auf der Konto-Sicherheitsseite einer Passkey-unterstuetzenden Website (z. B. Google, Microsoft, Amazon, eBay, PayPal) „Passkey hinzufuegen / Add passkey" auswaehlen.
  \item Das System zeigt ein Face-ID-/Touch-ID-Fenster; nach Bestaetigung wird der Passkey automatisch erstellt und im iCloud-Schluesselbund gespeichert.
  \item Beim naechsten Login auf der Website nach Eingabe des Benutzernamens „Mit iPhone anmelden" auswaehlen; Face ID einmal scannen – fertig.
  \item Verwaltung: Einstellungen Passwoerter entsprechende Website suchen; Passkey-Eintrag sichtbar.
\end{enumerate}

\textbf{Android-Pfad:}

\begin{enumerate}
  \item Auf der Konto-Sicherheitsseite einer Passkey-unterstuetzenden Website „Passkey hinzufuegen" auswaehlen.
  \item Google-Passwortmanager-Popup Google-Konto auswaehlen Bildschirmsperre verifizieren (Fingerabdruck/Muster/PIN).
  \item Passkey wird automatisch im Google-Passwortmanager gespeichert und zwischen Android-Geraeten desselben Kontos synchronisiert.
  \item Beim naechsten Login „Passkey verwenden" auswaehlen; Bildschirmsperre einmal entsperren – fertig.
  \item Verwaltung: Android-Einstellungen Passwoerter, Schluessel und Auto-Ausfuellung Passwortmanager google.com.
\end{enumerate}

\booksubsection{Passkey-Login ueber Geraetegrenzen hinweg}

Wenn du an einem Windows-PC eine Website aufrufst, fuer die nur dein iPhone einen Passkey hat, kannst du per QR-Code autorisieren:

\begin{enumerate}
  \item Auf der Login-Seite „Passkey verwenden" auswaehlen.
  \item „Anderes Geraet verwenden" auswaehlen.
  \item Der PC zeigt einen QR-Code; mit dem iPhone scannen.
  \item Das iPhone zeigt eine Face-ID-Bestaetigung; nach Bestaetigung ist der Login am PC abgeschlossen.
\end{enumerate}

Dieser Vorgang nutzt eine doppelte Verschluesselung per Bluetooth + Internet; der Private Key verlaesst niemals dein iPhone.

\begin{tcbpractice}{Jetzt ausprobieren: Passwortlosen Login erleben}
\begin{enumerate}
  \item Google-Konto-Login Sicherheit Passkeys Passkey hinzufuegen; Face-ID-/Fingerabdruck-Login erleben.
  \item Microsoft-Konto-Login Sicherheit Erweiterte Sicherheitsoptionen Passkey hinzufuegen.
  \item Bei Taobao, JD.com, eBay, PayPal und weiteren Passkey-unterstuetzenden Diensten ausprobieren.
  \item Traditionelles Passwort als Backup behalten; schrittweise hochwertige Konten auf Passkey migrieren.
\end{enumerate}

\end{tcbpractice}

\begin{tcbwarn}{Empfehlung fuer die Uebergangsphase}
Passkey hat Passwoerter noch nicht vollstaendig abgeloest – viele Websites sind im Uebergang „Passkey + Passwort“. Empfehlung: Hochwertige Konten zuerst mit Passkey aktivieren, aber\textbf{traditionelle Passwoerter nicht sofort loeschen}, als Backup behalten; weiterhin den Passwortmanager nutzen, da nicht alle Websites Passkey unterstuetzen. Passkey als „starke Ergaenzung” zum Passwortmanager betrachten, nicht als „sofortige Ersetzung“.

\end{tcbwarn}

\booksection{VII. Hardware-Sicherheitsschluessel}

Wenn deine Sicherheitsanforderungen ueber die normaler Nutzer hinausgehen – z. B. grosse Kryptowaehrungsbestaende, Vorstandsmitglied/Journalist/Aktivist oder Schutz von Unternehmenskonten – ist ein Hardware-Sicherheitsschluessel zuverlaessiger als eine SMS-Verifizierung.

\booksubsection{Was ist ein Hardware-Sicherheitsschluessel?}

Ein Hardware-Schluessel ist ein physisches Geraet (aehnlich einem USB-Stick) mit einem internen Sicherheitschip, der deinen Private Key speichert. Eine Beruehrung oder NFC ist fuer die Authentifizierung erforderlich. Massgebliche Produkte:

\begin{itemize}
  \item \textbf{YubiKey}(Yubico): Massgeblich; unterstuetzt FIDO2/WebAuthn, U2F, OTP, PIV, OpenPGP; USB-C / USB-A / NFC. Serie 5: 50-70 USD.
  \item \textbf{Google Titan}: Von Google; USB-C + NFC / Bluetooth; ca. 30-40 USD.
  \item \textbf{OnlyKey}: Open-Source-Alternative; unterstuetztmehrere Protokolle.
  \item \textbf{Chinesische Alternativen}: z. B. Feitian ePasskey; guenstiger.
\end{itemize}

\booksubsection{Warum sind Hardware-Schluessel sicherer als SMS-Codes?}

\begin{itemize}
  \item \textbf{Phishing-Schutz}: Wie Passkey antwortet der Hardware-Schluessel nur auf echte Domains; Phishing-Seiten koennen keinen Code erbeuten.
  \item \textbf{Schutz vor SIM-Swap}: SMS-Codes koennen per SIM-Swap abgefangen werden; der Hardware-Schluessel benoetigt keine Telefonnummer.
  \item \textbf{MITM-Schutz}: Das FIDO2-Protokoll enthaelt bei jeder Signatur eine Challenge; nicht wiederverwendbar.
  \item \textbf{Physische Isolierung}: Der Private Key bleibt immer im Hardware-Chip; selbst wenn Schadsoftware deinen Computer kompromittiert, ist er nicht lesbar.
\end{itemize}

\booksubsection{Einsatzbereiche}

\begin{itemize}
  \item Kryptoboesen-Konten, Wallet-Wiederherstellung.
  \item Unternehmens-SSO, Remote-VPN, SSH-Server-Login.
  \item E-Mail, GitHub, Social Media – „Master-Konten", deren Kompromittierung alle anderen Konten gefaehrden.
  \item Journalisten, Anwaelte, Führungskräfte und andere hochsensible Personen.
\end{itemize}

\booksubsection{Einrichtung eines Hardware-Schluessels als zweite Authentifizierung}

Am Beispiel eines Google-Kontos + YubiKey:

\begin{enumerate}
  \item YubiKey kaufen (empfohlen: zwei Stueck – einer fuer den Alltag, einer als Backup an unterschiedlichen Orten).
  \item Login auf myaccount.google.com Sicherheit Zwei-Schritt-Verifizierung.
  \item „Sicherheitsschluessel hinzufuegen" YubiKey in den USB-C-Port einfuegen oder an die NFC-Zone des Smartphones halten.
  \item Den Metallkontakt des YubiKey beruehren, um die Registrierung abzuschließen.
  \item Empfohlen: zweiten Backup-Schluessel registrieren.
  \item In den „Zwei-Schritt"-Einstellungen den Hardware-Schluessel als bevorzugte Methode festlegen; SMS-Code als Backup.
\end{enumerate}

Unterstützende Websites: Google, GitHub, Microsoft, Apple ID, Facebook, Twitter, Dropbox, GitLab, AWS, Coinbase, Binance u. a.

\begin{tcbinfo}{Zwei-Schluessel-Strategie}
Das groesste Risiko eines Hardware-Schluessels ist der Verlust. Empfehlung: 1. Zwei verschiedene Schluessel registrieren; 2. Einen dabei haben (Schluesselbund), einen im Haussafe oder Bankschliessfach; 3. Bei einigen Websites „Backup-Codes“ ausdrucken und offline aufbewahren. So bleibt man selbst bei Verlust nicht ausgesperrt.

\end{tcbinfo}

\begin{tcbwarn}{Kaufquelle sorgfaeltig waehlen}
Zwingend im offiziellen Shop oder bei autorisierten Haendlern kaufen; keine gebrauchten oder unabhaengigen Quellen – theoretisch koennte vorinstallierte Hardware einen Backdoor enthalten. YubiKey-Firmware ist nach Werk nicht aktualisierbar, aber Yubico veroeffentlicht fortlaufend Sicherheits-Hinweise auf der Website.

\end{tcbwarn}

\booksection{VIII. Passwort-Freigabe fuer Familie/Team}

In Familien werden Netflix, WLAN, Router-Admin und Smart-Lock-Passwoerter gemeinsam genutzt; in Teams sind es Server-, Datenbank- und Social-Media-Konten.\textbf{Passwoerter per WeChat zu senden oder in geteilten Dokumenten zu notieren ist hochriskant}, denn sobald Chatverlauf oder Dokumentberechtigungen geleakt werden, sind auch die Passwoerter kompromittiert.

\booksubsection{1Password „Shared Vault“}

1Password ist das Referenzprodukt fuer Familien- und Team-Sharing. Der Kernmechanismus ist\textbf{mehrere Vaults}:

\begin{itemize}
  \item \textbf{Persoenlicher Vault}: Nur fuer dich sichtbar; z. B. Bank, E-Mail und andere private Konten.
  \item \textbf{Shared Vault}: Fuer Familien- oder Teammitglieder sichtbar; z. B. Netflix, WLAN, gemeinsame Kreditkarte.
  \item \textbf{Berechtigungsstufen}: In der Team-Version koennen Mitglieder als „nur Lesen" oder „Lesen, aber nicht Bearbeiten" eingestellt werden.
\end{itemize}

Familien- oder Teammitglieder teilen nicht das Master-Passwort; jeder nutzt sein eigenes Konto, sieht aber die autorisierten Shared Vaults. Passwortaenderungen werden fuer alle in Echtzeit synchronisiert.

\booksubsection{Bitwarden „Organisation“-Funktion}

Die Gratis-Version von Bitwarden ermoeglicht das Teilen einzelner Passwoerter (per „Send“-Funktion zeitlich befristet); die Premium-Version kann „Organisationen” erstellen, um Eintraege mit Familienmitgliedern oder Kollegen zu teilen – mit feingranularen Berechtigungen. Bitwarden bietet zudem Self-Hosting fuer Unternehmen mit hohen Anforderungen an Datenhoheit.

\booksubsection{Familie: Netflix, WLAN und andere Passwoerter sicher teilen}

\begin{itemize}
  \item \textbf{Netflix / Video-Streaming}: Login-Passwort im „Shared Vault" von 1Password Families speichern; Familienmitglieder loggen sich selbst ein – kein Versenden im Gruppenchat.
  \item \textbf{WLAN-Passwort zu Hause}: iOS-Geraete koennen WLAN-Passwoerter per AirDrop direkt mit Gaesten teilen, ohne sie vorlesen; Android 12+ bietet „Wi-Fi per Nearby Share".
  \item \textbf{Smart-Lock, Smart-Home-Geraete}: Moeglichst die „Familienmitglieder"-Funktion der Hersteller-App nutzen und fuer jede Person ein eigenes Konto erstellen, statt einen Admin-Account zu teilen.
  \item \textbf{Router-Admin}: Admin-Passwort im 1Password Shared Vault speichern; fuer Gaeste ein separates SSID einrichten.
\end{itemize}

\begin{tcbinfo}{Preis-Leistungs-Verhaeltnis von 1Password Families}
1Password Families ca. 4,99 USD/Monat fuer bis zu 5 Familienmitglieder; jeder hat einen privaten Vault + 4 Shared Vaults. Pro Person weniger als 1 USD/Monat – die beste Familie-Sharing-Erfahrung. Bitwarden Familie: 10 USD/Jahr fuer 6 Personen; besseres Preis-Leistungs-Verhaeltnis, aber etwas schlechtere UI.

\end{tcbinfo}

\booksubsection{Team: Server-/Datenbank-Passwort-Sharing}

\begin{itemize}
  \item \textbf{1Password Teams / Business}: Mehrere Vaults pro Abteilung/Projektgruppe; feingranulare Berechtigungen; alle Zugriffe auditierbar.
  \item \textbf{Bitwarden Teams / Enterprise}: Open-Source; Self-Hosting moeglich; Enterprise-Version mit SSO-Integration und Verzeichnissynchronisierung.
  \item \textbf{Trennung der Zustaendigkeiten}: Admin-Passwoerter sollten nicht bei einer einzigen Person liegen; 2-3 „Notfallkontakte" benennen, damit bei Ausscheiden/Urlaub der Betrieb nicht stillsteht.
  \item \textbf{Austrittsprozess}: Beim Ausscheiden eines Mitarbeiters sofort dessen Passwortmanager-Zugriff entziehen; bei Bedarf alle von ihm beruehrten hochsensiblen Passwoerter rotieren.
\end{itemize}

\booksubsection{Notfallzugriff (Emergency Access): Falls dir etwas passiert}

Stell dir vor: Du verwaltest allein alle wichtigen Familienkonten; bei unerwartetem Krankenhausaufenthalt oder Ableben koennen Angehoerige nicht auf Bank, E-Mail oder Versicherungen zugreifen.\textbf{Die Funktion „Notfallzugriff“}ist genau dafuer gedacht – sie autorisiert eine vertrauenswuerdige Person, bei deiner Unerreichbarkeit auf deine Passwortdatenbank zuzugreifen.

\textbf{Bitwarden Notfallzugriff einrichten:}

\begin{enumerate}
  \item Bitwarden Web Vault Einstellungen Notfallzugriff.
  \item E-Mail-Adresse einer Vertrauensperson hinzufuegen und einladen.
  \item „Wartezeit" festlegen (empfohlen 7-14 Tage): Das Zeitfenster vom Antrag bis zu deiner Bestaetigung.
  \item Wenn du innerhalb der Wartezeit nicht reagierst (ablehnst oder nichts tust), erhaelt die Person nach Ablauf automatisch Zugriff.
  \item Du kannst jederzeit manuell ablehnen oder den Notfallzugriff widerrufen.
\end{enumerate}

\textbf{1Password Notfallzugriff:}

\begin{enumerate}
  \item 1Password nennt es nicht direkt „Notfallzugriff", aber ueber die „Familien-/Team-Organisation" realisierbar: Familienmitglieder hinzufuegen und bestimmte Eintraege im „Shared Vault" ablegen.
  \item Fuer Einzel-Nutzer: Das Master-Passwort einer hochgradig vertrauenswuerdigen Person mitteilen oder in einem versiegelten Umschlag im Safe verwahren.
  \item 1Password bietet zudem „Wiederherstellungscodes"; bei Aktivierung der Zwei-Faktor-Authentifizierung unbedingt ausdrucken und aufbewahren.
\end{enumerate}

\begin{tcbpractice}{Jetzt einrichten: Hinterlege einen „Reserveschluessel" fuer die Familie}
\begin{enumerate}
  \item Eine vertrauenswuerdige Person als „Notfallkontakt" bestimmen.
  \item Im Passwortmanager Notfallzugriff aktivieren (Bitwarden) oder dem Familien-Shared-Vault beitreten (1Password).
  \item Master-Passwort + Wiederherstellungscode auf Papier notieren, in Umschlag versiegeln, im Haussafe verwahren oder einer vertrauenswuerdigen Person mitteilen.
  \item Halbjaehrlich ueberpruefen, ob der Notfallkontakt noch passend ist.
\end{enumerate}

\end{tcbpractice}

\booksection{IX. Bewaehrte Praktiken fuer Passwortmanager}

\begin{itemize}
  \item \textbf{Starkes Master-Passwort festlegen}: Das Master-Passwort ist die letzte Verteidigungslinie – wählen Sie ein langes, komplexes Passwort und prägen Sie es sich ein.
  \item \textbf{Zwei-Faktor-Authentifizierung aktivieren}: Richten Sie für das Passwortmanager-Konto 2FA ein, für eine zusätzliche Sicherheitsebene.
  \item \textbf{Schwache Passwörter regelmäßig prüfen}: Die meisten Passwortmanager haben eine „Sicherheits-Audit"-Funktion, die schwache und doppelte Passwörter aufspürt.
  \item \textbf{Master-Passwort nicht vergessen}: Das Master-Passwort kann nicht wiederhergestellt werden – vergessen heißt, alle Passwörter zu verlieren. Am besten aufschreiben und an einem sicheren Ort aufbewahren.
  \item \textbf{Notfallzugang}: Richten Sie einen Notfallkontakt ein, damit eine Vertrauensperson im Notfall auf Ihre Passwörter zugreifen kann.
\end{itemize}

\begin{tcbwarn}{Wichtiger Hinweis}
Ein Passwortmanager ist praktisch, birgt aber auch Risiken: Wenn das Master-Passwort geleakt wird, sind alle Passwörter offen. Daher darf das Master-Passwort niemals mit anderen identisch sein und sollte nicht an leicht auffindbaren Orten notiert werden. Wählen Sie außerdem einen Anbieter mit guter Sicherheits-Historie.

\end{tcbwarn}

\begin{tcbpractice}{Jetzt ausprobieren: Starten Sie mit einem Passwortmanager}
\begin{enumerate}
  \item Wählen Sie einen Passwortmanager (Bitwarden oder systemeigen empfohlen).
  \item Registrieren Sie ein Konto und legen Sie ein starkes Master-Passwort fest.
  \item Importieren Sie vorhandene Passwörter (falls vorhanden) oder beginnen Sie, neue hinzuzufügen.
  \item Installieren Sie das Browser-Plugin oder die Smartphone-App und testen Sie das automatische Ausfüllen.
  \item Prüfen Sie den Sicherheits-Audit-Bericht und finden Sie Passwörter, die verbessert werden sollten.
\end{enumerate}

\end{tcbpractice}
